\documentclass[11pt]{article}

\usepackage[margin=1in]{geometry}
\usepackage{amsmath,amssymb,amsthm}
\usepackage{hyperref}

\newtheorem{theorem}{Theorem}[section]
\newtheorem{lemma}[theorem]{Lemma}
\newtheorem{proposition}[theorem]{Proposition}
\newtheorem{corollary}[theorem]{Corollary}
\theoremstyle{definition}
\newtheorem{definition}[theorem]{Definition}
\theoremstyle{remark}
\newtheorem{remark}[theorem]{Remark}

\newcommand{\HF}{\widehat{\mathrm{HF}}}
\newcommand{\C}{\mathbb{C}}
\newcommand{\Z}{\mathbb{Z}}
\newcommand{\Q}{\mathbb{Q}}
\newcommand{\R}{\mathbb{R}}

\title{Prime summands of surgeries on L-space knots}
\author{}
\date{}

\begin{document}

\maketitle

\begin{abstract}
We prove that every Dehn surgery on an L-space knot in \(S^3\) has at
most two nontrivial prime summands.  More precisely, a reducible
surgery occurs only at the cabling slope of a cable knot, and the two
summands are a lens space and the associated fractional surgery on the
companion.
\end{abstract}

\section{Statement}

We use the following standard convention.

\begin{definition}
An \emph{L-space} is a closed oriented rational homology sphere \(Y\)
such that
\[
\operatorname{rank}_{\Z/2}\HF(Y;\Z/2)
=
|H_1(Y;\Z)|.
\]
An \emph{L-space knot} is a nontrivial knot \(K\subset S^3\) admitting
a nonzero Dehn surgery which is an L-space.
\end{definition}

Some authors require the L-space surgery slope to be positive.  The
distinction is immaterial here: reversing orientation changes the sign
of a surgery slope and does not alter the number of prime summands.

The result is the following.

\begin{theorem}\label{thm:main}
Let \(K\) be an L-space knot in \(S^3\), and let
\[
Y=S^3_r(K)
\]
be obtained by Dehn surgery along \(K\).  Then \(Y\) has at most two
nontrivial prime summands.

More precisely:
\begin{enumerate}
\item if \(Y\) is irreducible, it has at most one nontrivial prime
      summand;
\item if \(Y\) is reducible, then \(K\) is a nontrivial cable
      \(C_{p,q}(J)\), the slope is the cabling slope \(pq\), and
      \[
      S^3_{pq}\bigl(C_{p,q}(J)\bigr)
      \cong
      L(p,q)\# S^3_{q/p}(J),
      \]
      where \(S^3_{q/p}(J)\) is irreducible and nontrivial whenever
      \(J\) is nontrivial.  Thus a reducible surgery has exactly two
      nontrivial prime summands.
\end{enumerate}
\end{theorem}

The unknot is usually excluded from the definition of an L-space knot,
but its surgeries are lens spaces, \(S^3\), or \(S^1\times S^2\), so
the same bound holds for the unknot as well.

\section{The reducible surgery theorem for L-space knots}

The key external classification input is the following established
special case of the cabling theorem.  It is important that this is not
the still-open general cabling conjecture: the L-space hypothesis is
essential.

\begin{theorem}[Reducible-surgery theorem for L-space knots]
\label{thm:red-lspace}
Let \(K\) be a nontrivial L-space knot.  If \(S^3_r(K)\) is
reducible, then \(K\) is a nontrivial cable and \(r\) is its cabling
slope.  After reversing orientations and adjusting cable conventions,
\[
K=C_{p,q}(J),\qquad r=pq,
\]
where \(p,q>1\) and \(\gcd(p,q)=1\).
\end{theorem}

\begin{proof}
This is the reducible-surgery specialization of the cabling theorem
for knots admitting L-space surgeries.  It follows from the
Floer-theoretic satellite classification of Hom--Lidman--Vafaee
\cite{HLV14} together with Krcatovich's reduced knot Floer criterion
for cabling \cite{Krc15}; the classical reducing-surface formulation
is treated in Scharlemann \cite{Scharlemann90}.  See also the
positivity and fiberedness results of Ozsv\'ath--Szab\'o, Ni, and
Hedden \cite{OS05,Ni07,Hed10}, which place L-space knots in the
class to which these results apply.
\end{proof}

We shall not use any general reducible-surgery statement for knots
that are not L-space knots.

\section{Cable surgery}

All slopes are measured with respect to the Seifert longitude and the
meridian.  For the standard \((p,q)\)-cable \(C_{p,q}(J)\), the cabling
slope is \(pq\).

\begin{lemma}[Cable-surgery decomposition]\label{lem:cable}
Let \(J\subset S^3\) be a knot, and let \(C_{p,q}(J)\) be its
\((p,q)\)-cable, with \(\gcd(p,q)=1\) and \(p>1\).  Then
\[
S^3_{pq}\bigl(C_{p,q}(J)\bigr)
\cong
L(p,q)\# S^3_{q/p}(J).
\]
The orientation convention for lens spaces may replace \(L(p,q)\) or
\(L(q,p)\) by its orientation reverse; this does not affect the
argument below.
\end{lemma}

\begin{proof}
Let \(V\) be a tubular neighborhood of \(J\), and let
\[
X=V\setminus \operatorname{int} N(C_{p,q})
\]
be the cable space.  The exterior of the cable knot is
\[
E(C_{p,q}(J))=E(J)\cup_{\partial V} X.
\]
Filling \(X\) along the cabling slope \(pq\) produces a reducing
\(2\)-sphere.  One side of this sphere is a punctured lens space
\(L(p,q)\), while the other side is the companion exterior \(E(J)\)
filled along the induced slope \(q/p\).  Hence the filled manifold is
the claimed connected sum.  This computation is classical; see Gordon
\cite{Gordon83} or Scharlemann \cite{Scharlemann90}.

For example, when \(J\) is the unknot, the formula becomes
\[
S^3_{pq}(T(p,q))
\cong
L(p,q)\# L(q,p),
\]
which also checks the homology: both sides have first homology of
order \(pq\).
\end{proof}

The following theorem of Gordon and Luecke will be used to show that
the companion-surgery factor is prime.

\begin{theorem}[Gordon--Luecke]\label{thm:integral}
If a nontrivial Dehn surgery on a nontrivial knot in \(S^3\) is
reducible, then the surgery slope is integral.
\end{theorem}

\begin{proof}
See Gordon--Luecke \cite{GL87}, where it is proved that only integral
Dehn surgeries on nontrivial knots in \(S^3\) can yield reducible
manifolds.
\end{proof}

\begin{lemma}\label{lem:companion-prime}
Let \(p,q>1\) be coprime, and let \(J\) be nontrivial.  Then
\(S^3_{q/p}(J)\) is a nontrivial prime rational homology sphere.
\end{lemma}

\begin{proof}
Since \(\gcd(p,q)=1\), homological calculation gives
\[
H_1\bigl(S^3_{q/p}(J);\Z\bigr)\cong \Z/q\Z.
\]
Because \(q>1\), the manifold is not \(S^3\).

If \(S^3_{q/p}(J)\) were reducible, Theorem \ref{thm:integral} would
force the slope \(q/p\) to be integral.  But \(q/p\) is reduced and
\(p>1\), so it is not integral.  Thus \(S^3_{q/p}(J)\) is
irreducible.  Being irreducible and not \(S^3\), it is prime.
\end{proof}

\section{Zero surgery}

The slope \(0\) is not covered by the rational homology sphere
argument.  For L-space knots it can be treated directly.

\begin{theorem}[Ozsv\'ath--Szab\'o and Ni]\label{thm:fibered}
Every L-space knot is fibered.
\end{theorem}

\begin{proof}
Ozsv\'ath and Szab\'o proved that the Alexander polynomial of an
L-space knot is monic \cite{OS05}.  Ni's fibering theorem
\cite{Ni07} says that knot Floer homology detects fiberedness;
together these imply that an L-space knot is fibered.
\end{proof}

\begin{lemma}\label{lem:zero}
If \(K\) is a nontrivial L-space knot, then \(S^3_0(K)\) is
irreducible.
\end{lemma}

\begin{proof}
By Theorem \ref{thm:fibered}, \(K\) is fibered.  Let \(F\) be its
fiber surface, of genus \(g\).  Since \(K\) is nontrivial, \(g\geq 1\).

In \(0\)-surgery, the boundary of a page is capped by the meridian
disk of the surgery solid torus.  Consequently \(S^3_0(K)\) fibers
over \(S^1\) with fiber the closed genus-\(g\) surface
\(\overline F\).  Thus \(S^3_0(K)\) is a mapping torus
\[
S^3_0(K)\cong
(\overline F\times [0,1])/(x,1)\sim(\phi(x),0).
\]
The universal cover of this mapping torus is
\[
\widetilde{\overline F}\times\R\cong \R^2\times\R\cong\R^3,
\]
because \(g\geq 1\).  Hence the mapping torus is irreducible.
\end{proof}

\section{Proof of the main theorem}

\begin{proof}[Proof of Theorem \ref{thm:main}]
If \(K\) is the unknot, then every surgery is \(S^3\), a lens space,
or \(S^1\times S^2\), so there is at most one nontrivial prime
summand.  Henceforth assume that \(K\) is nontrivial.

If \(r=\infty\), then \(Y=S^3\), which has no nontrivial prime
summands.  If \(r=0\), Lemma \ref{lem:zero} shows that \(Y\) is
irreducible.

Now suppose \(r\) is finite and nonzero.  If \(Y\) is irreducible,
there is at most one nontrivial prime summand.  Suppose therefore
that \(Y\) is reducible.

By Theorem \ref{thm:red-lspace}, after reversing orientations if
necessary,
\[
K=C_{p,q}(J),\qquad r=pq,
\]
with \(p,q>1\) and \(\gcd(p,q)=1\).  Lemma \ref{lem:cable} then gives
\[
Y\cong L(p,q)\# S^3_{q/p}(J).
\]

The first factor is a nontrivial lens space because \(p>1\), and lens
spaces are prime.  If \(J\) is the unknot, the second factor is the
nontrivial lens space
\[
S^3_{q/p}(U)\cong L(q,p),
\]
since \(q>1\).  If \(J\) is nontrivial, Lemma
\ref{lem:companion-prime} says that the second factor is nontrivial
and prime.  Thus \(Y\) has exactly two nontrivial prime summands.

By the Kneser--Milnor prime decomposition theorem \cite{Milnor62},
these are the prime summands of \(Y\).  Therefore every Dehn surgery
on an L-space knot has at most two prime summands.
\end{proof}

\section{The L-space case}

If one begins with a surgery that is itself an L-space, the two
summands are also L-spaces.

\begin{proposition}\label{prop:summands}
Every prime summand of a connected sum of L-spaces is an L-space.
\end{proposition}

\begin{proof}
Let \(Y=Y_1\#\cdots\#Y_n\).  The connected-sum K\"unneth theorem for
\(\HF\) gives
\[
\operatorname{rank}\HF(Y)
=
\prod_{i=1}^n \operatorname{rank}\HF(Y_i),
\]
while
\[
|H_1(Y;\Z)|
=
\prod_{i=1}^n |H_1(Y_i;\Z)|.
\]
For every rational homology sphere,
\[
\operatorname{rank}\HF(Y_i)\geq |H_1(Y_i;\Z)|.
\]
If \(Y\) is an L-space, equality must hold for every factor.  Thus
every \(Y_i\) is an L-space.
\end{proof}

\begin{corollary}
If an L-space obtained by surgery on an L-space knot is reducible,
then it is a connected sum of exactly two prime L-spaces.  In the
notation of Theorem \ref{thm:main}, one summand is the lens space
\(L(p,q)\), and the other is \(S^3_{q/p}(J)\).
\end{corollary}

\begin{proof}
This is immediate from Theorem \ref{thm:main} and Proposition
\ref{prop:summands}.
\end{proof}

\begin{thebibliography}{99}

\bibitem[CGLS87]{CGLS87}
M.~Culler, C.~McA.~Gordon, J.~Luecke, and P.~B.~Shalen,
\emph{Dehn surgery on knots},
Ann. of Math. (2) \textbf{125} (1987), 237--300.

\bibitem[Gor83]{Gordon83}
C.~McA.~Gordon,
\emph{Dehn surgery and satellite knots},
Trans. Amer. Math. Soc. \textbf{275} (1983), 687--708.

\bibitem[GL87]{GL87}
C.~McA.~Gordon and J.~Luecke,
\emph{Only integral Dehn surgeries can yield reducible manifolds},
Math. Proc. Cambridge Philos. Soc. \textbf{102} (1987), 97--101.

\bibitem[Hed10]{Hed10}
M.~Hedden,
\emph{Notions of positivity and the Ozsv\'ath--Szab\'o concordance
invariant},
J. Knot Theory Ramifications \textbf{19} (2010), 617--629.

\bibitem[HLV14]{HLV14}
J.~Hom, T.~Lidman, and F.~Vafaee,
\emph{Berge--Gabai knots and L-space satellite operations},
Algebr. Geom. Topol. \textbf{14} (2014), 3745--3763.

\bibitem[Krc15]{Krc15}
D.~Krcatovich,
\emph{The reduced knot Floer complex},
Topology Appl. \textbf{194} (2015), 171--200.

\bibitem[Mil62]{Milnor62}
J.~Milnor,
\emph{A unique decomposition theorem for 3-manifolds},
Amer. J. Math. \textbf{84} (1962), 1--7.

\bibitem[Ni07]{Ni07}
Y.~Ni,
\emph{Knot Floer homology detects fibred knots},
Invent. Math. \textbf{170} (2007), 577--608.

\bibitem[OS05]{OS05}
P.~Ozsv\'ath and Z.~Szab\'o,
\emph{On knot Floer homology and lens space surgeries},
Topology \textbf{44} (2005), 1281--1300.

\bibitem[Sch90]{Scharlemann90}
M.~Scharlemann,
\emph{Producing reducible 3-manifolds by surgery on a knot},
Topology \textbf{29} (1990), 481--500.

\end{thebibliography}

\end{document}
